\documentclass[12pt,a4paper]{article}

% Documento final de Mitotl IA. La información entre corchetes se completa
% durante las versiones V1-V4.
\usepackage[utf8]{inputenc}
\usepackage[T1]{fontenc}
\usepackage{lmodern}
\usepackage{geometry}
\usepackage{portada}
\usepackage{setspace}
\usepackage{microtype}
\usepackage{graphicx}
\usepackage{booktabs}
\usepackage{tabularx}
\usepackage{array}
\usepackage{amsmath}
\usepackage{xcolor}
\usepackage{tikz}
\usetikzlibrary{arrows.meta,positioning,shapes.geometric}
\usepackage{fancyhdr}
\usepackage{enumitem}
\usepackage{csquotes}
\usepackage[backend=biber,style=apa]{biblatex}
\usepackage[spanish,es-nodecimaldot,es-noquoting]{babel}
\usepackage[hidelinks]{hyperref}
\addbibresource{referencias.bib}

\geometry{margin=2.5cm}
\onehalfspacing
\setlength{\parindent}{0pt}
\setlength{\parskip}{0.65em}

\definecolor{mitotlblue}{HTML}{064653}
\definecolor{mitotlteal}{HTML}{087A5A}
\definecolor{mitotlpink}{HTML}{FF39B0}
\definecolor{placeholder}{HTML}{6B7280}

\hypersetup{
  pdftitle={Mitotl IA: análisis y retroalimentación de movimiento corporal},
  pdfauthor={David Magaña Celis},
  colorlinks=true,
  linkcolor=mitotlblue,
  citecolor=mitotlblue,
  urlcolor=mitotlblue
}

\pagestyle{fancy}
\fancyhf{}
\lhead{Mitotl IA}
\rhead{}
\cfoot{\thepage}
\setlength{\headheight}{15pt}


\newcommand{\autor}{[Nombre del autor]}
\newcommand{\institucion}{Facultad de Estudios Superiores Acatlán}
\newcommand{\curso}{Módulo 5 --- Diplomado en Ciencia de Datos}
\newcommand{\fechaentrega}{05 de agosto de 2026}
\newcommand{\pendiente}[1]{\textcolor{placeholder}{\textit{Pendiente V1--V4: #1}}}

\title{Análisis, comparación y retroalimentación de movimiento corporal en videos de danza}
\author{David Magaña Celis}
\faculty{\institucion}
\degree{\curso}
\supervisor{Mtro. Fernando Barranco Rodríguez}
\cityandyear{\fechaentrega}
\logouni{figuras/UNAM.png}
\logofac{figuras/FESACATLAN.png}

\begin{document}

\maketitle

\tableofcontents
\clearpage

\section*{Resumen ejecutivo}
\addcontentsline{toc}{section}{Resumen ejecutivo}
Mitotl IA es una herramienta de apoyo para la práctica autónoma de danza. A partir de un video de referencia y un video de ejecución, el sistema identifica diferencias observables en el movimiento corporal y las convierte en una retroalimentación orientada a la práctica. El objetivo no es emitir un juicio artístico profesional, sino ayudar a responder preguntas concretas: en qué momento aparece una diferencia, qué parte del cuerpo conviene revisar y qué acción puede intentar la persona usuaria.

La versión actual trabaja con una sola persona y requiere que el cuerpo completo sea visible. Su flujo combina validación de video, extracción de pose, variables geométricas, sincronización temporal y reglas de retroalimentación. El resultado se presenta mediante un score orientativo, indicadores por parte del cuerpo, visualizaciones comparativas, clips críticos y un asistente conversacional básico. La validación se realizó con un video de referencia y una ejecución propia; por tanto, no constituye una evaluación general de desempeño artístico.

\section{Introducción}
Aprender una coreografía mediante videos permite practicar en horarios y espacios diversos, pero también desplaza a la persona practicante la tarea de detectar sus propios errores. La comparación visual directa puede ser insuficiente cuando la diferencia ocurre durante un instante, afecta una articulación específica o se presenta con una velocidad distinta entre la referencia y la ejecución.

Este documento presenta el diseño y desarrollo de Mitotl IA como una solución aplicada de ciencia de datos. Se describe el problema de uso, la estrategia de la versión actual, la procedencia de los datos, el procesamiento previsto y la evidencia de la aplicación. El reporte distingue entre lo que está implementado y las capacidades que permanecen como trabajo futuro.

\section{Planteamiento del problema}
\subsection{Nicho y contexto de uso}
El nicho de uso es el aprendizaje autónomo de danza mediante videos de referencia. Incluye a personas que practican coreografías, ejercicios o tendencias difundidas en plataformas como TikTok y YouTube, así como a quienes siguen una clase grabada sin contar necesariamente con un espejo, un salón especializado o la observación continua de un instructor.

En este contexto, el video de referencia funciona como una guía visual y el video de ejecución como evidencia de la práctica. La necesidad no es ordenar a las personas por nivel ni establecer una calificación universal de la danza, sino ofrecer información localizada que permita repetir y ajustar una secuencia.
\subsection{Usuario objetivo}
La persona usuaria objetivo practica de manera individual y puede tener experiencia inicial o intermedia. Cuenta con un teléfono o una computadora para cargar o grabar un video, desea aprender una secuencia concreta y necesita una explicación comprensible sobre sus principales oportunidades de mejora.

Se asume que la persona puede colocar la cámara de frente y mantener el cuerpo completo dentro del encuadre. Este supuesto es importante porque la versión actual no está diseñada para grupos, personas parcialmente visibles ni situaciones en las que una parte importante del cuerpo quede tapada o fuera de la vista.
\subsection{Problema y necesidad}
El problema es que una persona puede repetir una coreografía sin saber con precisión qué está fallando. La observación subjetiva permite reconocer que algo no coincide, pero no siempre permite localizar el momento, distinguir si la diferencia proviene de brazos, piernas o torso, ni convertir esa observación en una instrucción práctica.

De este problema se desprenden cuatro necesidades: localizar momentos específicos de dificultad; identificar las partes del cuerpo asociadas; recibir una explicación clara y accionable; y contar con una guía que permita volver a practicar la sección relevante. Mitotl IA atiende estas necesidades mediante comparación corporal y temporal, sin presentarse como sustituto de una evaluación profesional.
\subsection{Propuesta de valor y caso de uso}
La propuesta de valor de Mitotl IA es transformar dos videos en una guía de práctica localizada. La persona carga un video de referencia y un video de ejecución, espera el análisis y consulta un resultado compuesto por similitud corporal, similitud temporal, hallazgos y recomendaciones. También puede revisar una visualización comparativa y preguntar al asistente sobre la interpretación del resultado.

El caso de uso principal es una sesión de práctica individual: seleccionar una secuencia, grabar o cargar la ejecución, revisar las diferencias más relevantes y repetir la sección con base en la retroalimentación. El sistema no evalúa calidad artística profesional, musicalidad, expresividad o intención escénica. Su salida es una aproximación computacional orientada a comparar movimiento observable.

\section{Estrategia de solución}
\subsection{Accionables}
La estrategia prioriza resultados que puedan convertirse en acciones inmediatas. A partir del análisis, la persona puede: revisar el segmento temporal con mayor diferencia; concentrarse en una parte del cuerpo; observar la referencia y la ejecución en paralelo; repetir la sección crítica; y utilizar el asistente para aclarar el significado de un indicador.

El score general se comunica como una señal de orientación y no como una calificación definitiva. La interfaz debe favorecer la interpretación por evidencia: score, momento, parte del cuerpo y recomendación se presentan juntos para evitar que un número aislado sea la única conclusión.
\subsection{Usabilidad y usuarios}
El flujo está organizado alrededor de cuatro momentos: cargar los videos, ejecutar el análisis, revisar resultados y consultar la retroalimentación. La interfaz conserva una navegación por pestañas para separar videos, resultados, visualizaciones y asistente. Para una demostración, también se dispone de una sesión precalculada que permite revisar la interfaz sin esperar nuevamente todo el procesamiento.

Los mensajes de la aplicación distinguen advertencias de la referencia y de la ejecución. Esto es relevante cuando un video contiene frames que no pudieron decodificarse o momentos sin detección corporal: la advertencia informa una condición técnica sin convertirla automáticamente en un juicio sobre la persona.
\subsection{Alcance de la versión actual}
La versión actual incluye análisis de videos cargados, extracción de pose con MediaPipe, comparación de movimiento, alineación temporal, score orientativo, recomendaciones, visualización comparativa, clips críticos y un asistente básico. La captura mediante cámara y el análisis posterior forman parte del flujo previsto; el feedback perfectamente en tiempo real no es un requisito de esta versión.

Quedan fuera del alcance inmediato la evaluación artística profesional, el análisis de varias personas, la personalización longitudinal, los planes de práctica adaptativos, la sincronización musical avanzada y los modelos especializados entrenados con etiquetas de calidad de danza. Estas capacidades pueden explorarse en versiones futuras después de reunir datos autorizados y criterios de evaluación adecuados.

\section{Descripción de los datos}
\subsection{Fuentes de referencia y ejecución}
La fuente académica principal para el video de referencia es AIST Dance Video Database, una colección de videos de danza utilizada para construir y validar el flujo inicial. Se selecciona una secuencia con una persona, vista frontal y cuerpo completo visible. La referencia se emplea para probar la extracción de pose, la comparación corporal y la sincronización temporal, no para definir una verdad absoluta sobre cómo debe bailar una persona.

El video de ejecución representa la práctica propia que se desea comparar. Para validar el flujo se utilizó una ejecución grabada por la persona desarrolladora, distinta del video de referencia. Esta evidencia permite revisar el funcionamiento de la comparación, pero no basta para medir aprendizaje, generalización o calidad de danza en una población amplia.

Como fuente complementaria se considera AIST++, que aporta anotaciones y referencias de movimiento útiles para explorar landmarks, confianza y cobertura corporal. Su función es apoyar la validación técnica; no se utiliza como conjunto etiquetado de personas que bailan bien o mal.
\subsection{Restricciones de uso}
El uso de los videos de AIST está sujeto a sus términos de uso y a las condiciones académicas de la fuente \cite{aist,aistterms}. Los archivos de video utilizados localmente no deben redistribuirse dentro del repositorio. Por esa razón, el proyecto conserva metadatos y rutas esperadas, pero excluye los videos privados o descargados de los artefactos versionados.

Los videos aportados por una persona usuaria deberán contar con autorización o una licencia compatible con el análisis. La aplicación debe tratar estos archivos como entradas temporales de procesamiento y evitar presentarlos como material público. La documentación separa explícitamente procedencia, permiso de uso y datos derivados para que una ejecución reproducible no dependa de un archivo privado.
\subsection{Datos derivados}
Mitotl IA genera datos derivados a partir de cada video. Entre ellos se encuentran los landmarks corporales, coordenadas normalizadas, visibilidad de puntos, ángulos articulares, trayectorias, marcas de tiempo, segmentos temporales, diferencias entre referencia y ejecución y scores. Estos datos no son una nueva fuente externa: son resultados calculados por el pipeline.

La representación derivada permite comparar movimiento sin depender directamente de la resolución o del tamaño original del video. También permite producir evidencia comprensible para la interfaz, como diferencias por parte del cuerpo y momentos críticos. En etapas posteriores se documentarán con mayor detalle las variables, la normalización, la modelación y las reglas de retroalimentación.

\section{Análisis exploratorio y calidad de entrada}
\subsection{Calidad de los videos}
Antes de analizar el movimiento, el sistema inspecciona la extensión, la duración, el FPS, la resolución y la cantidad de frames declarada por el contenedor. Después genera una copia temporal de análisis con una altura máxima de 480 píxeles y un FPS objetivo de 30. Esta preparación reduce el costo computacional y conserva el archivo original sin modificarlo.

La validación distingue entre errores que impiden continuar y advertencias que permiten generar resultados parciales. Por ejemplo, si el video reporta más frames de los que finalmente pueden decodificarse, se informa la diferencia y se continúa con los frames disponibles. La advertencia es relevante para interpretar el final del video, pero no implica que todo el análisis deba descartarse.
\subsection{Detección corporal}
MediaPipe Pose extrae hasta 33 landmarks por frame. Para cada punto se conservan sus coordenadas $x$, $y$, $z$ y un valor de visibilidad. La exploración revisa si existe detección, si los landmarks tienen la cantidad esperada y si los índices de frame y timestamps mantienen continuidad.

El análisis se limita a una persona con el cuerpo completo visible y una cámara frontal o aproximadamente frontal. Cuando un frame no contiene detección, se conserva su posición temporal para no desplazar artificialmente la secuencia. Los huecos breves pueden tratarse durante la construcción de la secuencia angular; los huecos largos se reportan como una condición que reduce la confiabilidad.
\subsection{Visualizaciones exploratorias}
Las visualizaciones exploratorias superponen landmarks sobre frames seleccionados y contrastan poses de referencia y ejecución. También permiten inspeccionar el comportamiento de brazos, piernas, torso y cabeza, además de ubicar los segmentos con mayor diferencia. Estas vistas sirven como control de calidad: antes de interpretar un score se verifica que la detección corresponda realmente a la persona y al movimiento observado.

\section{Ingeniería de variables}
\subsection{Extracción y normalización}
La extracción conserva coordenadas y marcas de tiempo por frame. Para hacer comparables dos videos con distinta escala o distancia a la cámara, las coordenadas se trasladan usando el centro de las caderas como origen y se escalan con la distancia entre el centro de hombros y el centro de caderas. Así, la comparación se concentra en la configuración relativa del cuerpo y no únicamente en su posición dentro de la imagen.

Además de las coordenadas normalizadas, se calculan ocho ángulos articulares: codos, hombros, rodillas y caderas de ambos lados. La visibilidad se conserva como diagnóstico de calidad. También se derivan velocidades angulares y velocidades de landmarks cuando existen timestamps válidos, aunque el score principal utiliza la información corporal y temporal definida para esta versión.
\subsection{Variables de comparación}
La comparación alineada calcula diferencias euclidianas en el plano visible $XY$ y una diferencia espacial auxiliar en $XYZ$ para cada landmark compartido. Las diferencias se resumen por frame, landmark y parte corporal. Las partes amplias usadas en los resultados son brazos, piernas, torso y cabeza; los brazos y piernas se obtienen agrupando sus lados izquierdo y derecho.

La secuencia de ocho ángulos se utiliza para sincronizar referencia y ejecución. Las variables temporales incluyen pares de frames alineados, desplazamiento de tiempo, duración de cada segmento, razón entre duraciones y similitud temporal. La coreografía se divide en segmentos de aproximadamente dos segundos para que los hallazgos puedan relacionarse con momentos revisables por la persona usuaria.
\subsection{Preparación para visualización y feedback}
Una diferencia corporal no se presenta directamente como una lista de coordenadas. Se transforma en una similitud acotada mediante $s(d)=1/(1+d)$, se agrupa por parte del cuerpo y se convierte a una escala de 0 a 100. Los segmentos débiles se ordenan por sus diferencias y se acompañan con tiempos de referencia y ejecución.

La estructura de resultados alimenta tanto las tarjetas y tablas de la interfaz como las recomendaciones y el contexto del asistente. Si falta información suficiente, el pipeline conserva el error o la advertencia por etapa para que la interfaz comunique qué ocurrió y si el resultado puede interpretarse con cautela.

\section{Modelación}
\subsection{Enfoque híbrido de la versión actual}
La solución combina componentes deterministas y un modelo preentrenado. MediaPipe Pose estima landmarks; las variables normalizadas y angulares describen el movimiento; la comparación geométrica calcula diferencias; DTW alinea secuencias realizadas a velocidades distintas; y las reglas convierten los resultados en retroalimentación. Este enfoque es adecuado para el alcance actual porque permite explicar cómo se obtiene cada indicador y no requiere entrenar un clasificador de calidad artística.
\subsection{Sincronización temporal con DTW}
Dynamic Time Warping (DTW) encuentra un camino de correspondencias entre la secuencia angular de referencia y la de ejecución. Antes de calcularlo, cada uno de los ocho ángulos se estandariza usando la media y desviación combinadas de ambas secuencias. Esto reduce el efecto de escalas distintas entre variables y permite comparar la forma temporal del movimiento.

La implementación conserva el camino DTW y lo traduce a pares de frames y timestamps. Los frames sin pose o con ángulos incompletos se mantienen en la línea temporal; las ausencias breves se interpolan y un hueco mayor a dos segundos se rechaza para evitar fabricar una trayectoria corporal poco confiable. El camino se proyecta después en segmentos de dos segundos.
\subsection{Score y reglas de retroalimentación}
El score general combina similitud corporal y similitud temporal con pesos explícitos:

\[
\text{score general}=100\left(0.80\,\text{similitud corporal}+0.20\,\text{similitud temporal}\right).
\]

La similitud corporal se calcula como el promedio de los grupos brazos, piernas, torso y cabeza. La similitud temporal se obtiene de la alineación y de la coincidencia de duración de los segmentos. El resultado se comunica como preliminar y orientativo, no como porcentaje de talento, error absoluto o calidad profesional.

Las reglas priorizan los hallazgos con mayor diferencia y los traducen en mensajes de práctica: revisar una parte corporal, volver a un segmento temporal o repetir la secuencia observando una articulación. El asistente explica estos resultados sin inventar métricas que no estén disponibles.

\section{Pipeline e inferencia}
\begin{figure}[htbp]
  \centering
  \begin{tikzpicture}[
    node distance=0.55cm and 0.35cm,
    boxstep/.style={draw=mitotlblue, fill=mitotlblue!8, rounded corners, align=center, minimum height=0.75cm, text width=2.4cm, font=\small},
    result/.style={draw=mitotlteal, fill=mitotlteal!10, rounded corners, align=center, minimum height=0.75cm, text width=2.4cm, font=\small},
    arrow/.style={-{Latex[length=2mm]}, thick, draw=mitotlblue}
  ]
    \node[boxstep] (input) {Videos de referencia\\y ejecución};
    \node[boxstep, right=of input] (validate) {Validación\\y preparación};
    \node[boxstep, right=of validate] (pose) {Extracción\\de pose};
    \node[boxstep, below=of pose] (features) {Variables\\corporales};
    \node[boxstep, left=of features] (dtw) {DTW y\\segmentación};
    \node[boxstep, left=of dtw] (compare) {Comparación\\corporal};
    \node[result, below=of compare] (score) {Score, hallazgos\\y recomendaciones};
    \node[result, right=of score] (ui) {Interfaz, clips\\y agente};
    \draw[arrow] (input) -- (validate);
    \draw[arrow] (validate) -- (pose);
    \draw[arrow] (pose) -- (features);
    \draw[arrow] (features) -- (dtw);
    \draw[arrow] (dtw) -- (compare);
    \draw[arrow] (compare) -- (score);
    \draw[arrow] (score) -- (ui);
  \end{tikzpicture}
  \caption{Flujo general de inferencia de Mitotl IA.}
  \label{fig:pipeline}
\end{figure}
El pipeline recibe dos videos y produce una sesión estructurada. Primero valida y prepara cada entrada; después extrae pose y variables; construye la alineación temporal; compara landmarks en cada par DTW; calcula scores y segmentos; genera feedback; y entrega resultados a la interfaz, los clips y el asistente.

Los errores se controlan por etapa. Una entrada inválida detiene el análisis con un mensaje específico; una copia temporal incompleta conserva los frames decodificados y advierte sobre el tramo omitido; un frame sin detección no se confunde con un frame inexistente; y un hueco largo de pose evita que DTW continúe con datos artificiales. Esta separación permite distinguir un problema del archivo de un resultado bajo de similitud.

\section{Resultados}
\subsection{Resultados cuantitativos}
\pendiente{Score general, scores corporales, similitud temporal y segmentos débiles.}
\subsection{Hallazgos y recomendaciones}
\pendiente{Partes del cuerpo, momentos críticos y recomendaciones generadas.}
\subsection{Visualización comparativa}
\pendiente{Video comparativo y clips críticos de la demo.}

\section{Interfaz de uso}
\subsection{Flujo de análisis}
\pendiente{Carga, ejecución, consulta de resultados y asistencia.}
\subsection{Demo precalculada}
\pendiente{Sesión guardada para revisar la interfaz sin esperar el análisis.}

\section{Evaluación y reproducibilidad}
\subsection{Pruebas del pipeline}
\pendiente{Pruebas unitarias, validación de una ejecución propia y validación visual.}
\subsection{Despliegue}
\pendiente{Ejecución local, Streamlit Cloud y dependencias.}
\subsection{Manejo de errores}
\pendiente{Frames sin detección, videos truncados, MediaPipe y ausencia de API key.}

\section{Limitaciones, ética y uso responsable}
\pendiente{Alcance de una persona, score orientativo, calidad de captura, licencias y privacidad.}

\section{Conclusiones y trabajo futuro}
\subsection{Conclusiones}
\pendiente{Hallazgos principales y respuesta al problema.}
\subsection{Trabajo futuro}
\pendiente{Tiempo real, varias personas, historial, personalización y métricas de ritmo.}

\section*{Referencias}
\addcontentsline{toc}{section}{Referencias}
\nocite{aist,aistterms,mediapipe,dtw}
\printbibliography[heading=none]

\end{document}
